% BOZZA (passo 2 del piano): bilancio dei parametri di una proiezione STERA e risparmio al crescere del modello.
% Destinazione proposta: nuova sezione del Capitolo 4 (formalismo), dopo la sezione sul gradiente.
% Non ancora inclusa in main.tex.

\section{Bilancio dei parametri}
\label{sec:bilancio-parametri}

Una proiezione di STERA non è un singolo blocco RowMask23, ma un mosaico di blocchi a cui si aggiungono tre componenti più piccole. Questa sezione conta esattamente i valori memorizzati e ricava il risparmio rispetto alla forma densa, prima per una proiezione, poi per uno strato del Transformer e infine per il modello intero.

\begin{definizione}[Proiezione STERA]
\label{def:proiezione-stera}
Siano $m$ e $n$ multipli di 128. Una proiezione STERA $W\in\mathbb{R}^{m\times n}$ è
\[
W \;=\; \operatorname{diag}(g_{\mathrm{out}})
\begin{bmatrix}
s_{11}\,\Phi(c_{11}) & \cdots & s_{1N}\,\Phi(c_{1N})\\
\vdots & \ddots & \vdots\\
s_{M1}\,\Phi(c_{M1}) & \cdots & s_{MN}\,\Phi(c_{MN})
\end{bmatrix}
\operatorname{diag}(g_{\mathrm{in}}) \;+\; U V^{\top},
\]
con $M=m/128$ e $N=n/128$ blocchi per lato. Le sue componenti sono:
\begin{itemize}
\item le \emph{coordinate} $c_{pq}\in\mathbb{R}^{3176}$ di ogni blocco RowMask23;
\item una \emph{scala} $s_{pq}\in\mathbb{R}$ per blocco;
\item i \emph{guadagni} per riga e per colonna $g_{\mathrm{out}}\in\mathbb{R}^{m}$ e $g_{\mathrm{in}}\in\mathbb{R}^{n}$;
\item il \emph{termine a basso rango} $UV^{\top}$, con $U\in\mathbb{R}^{m\times r}$ e $V\in\mathbb{R}^{n\times r}$.
\end{itemize}
\end{definizione}

\begin{proposizione}[Valori memorizzati e risparmio di una proiezione]
\label{prop:risparmio-proiezione}
La proiezione della Definizione~\ref{def:proiezione-stera} memorizza
\[
P(m,n,r) \;=\; \frac{3177}{16384}\,mn \;+\; (1+r)\,(m+n)
\]
valori. Il risparmio rispetto ai $mn$ valori della forma densa è
\[
s(m,n,r) \;=\; 1-\frac{3177}{16384}-(1+r)\left(\frac1m+\frac1n\right)
\;\approx\; 80{,}61\% \;-\; (1+r)\left(\frac1m+\frac1n\right).
\]
\end{proposizione}

\begin{proof}
I blocchi sono $mn/16384$. Ciascuno contribuisce con 3176 coordinate e una scala, cioè $3177\,mn/16384$ valori in tutto. I guadagni aggiungono $m+n$ valori e i fattori $U$ e $V$ ne aggiungono $r(m+n)$. Dividendo per $mn$ si ottiene $s$.
\end{proof}

La formula separa ciò che il formato garantisce da ciò che dipende dalle scelte di progetto:
\begin{itemize}
\item \textbf{coordinate e scale} occupano sempre la frazione $3177/16384=19{,}39\%$ della forma densa, qualunque siano le dimensioni: è il risparmio dell'$80{,}61\%$ garantito da RowMask23;
\item \textbf{i guadagni} dipendono solo dalla forma della matrice e pesano $1/m+1/n$, al più l'$1{,}56\%$ per una matrice $128\times128$;
\item \textbf{il basso rango} dipende dalla scelta di $r$ ed è l'unica componente che può spostare sensibilmente il risparmio.
\end{itemize}

\begin{osservazione}[Limiti]
Per ogni scelta dei parametri $s(m,n,r)<80{,}61\%$. In assenza del termine a basso rango ($r=0$), e poiché $m,n\ge128$, vale $s(m,n,0)\ge 80{,}61\%-\tfrac{2}{128}=79{,}05\%$.
\end{osservazione}

\paragraph{Uno strato del Transformer.} In uno strato con dimensione $d$ e feed-forward $f=4d$, le proiezioni sono quattro matrici $d\times d$ (attenzione) e due matrici $4d\times d$ e $d\times 4d$ (feed-forward), per $12d^2$ valori densi. Se il rango è proporzionale alla dimensione, $r=\kappa d$, sommando la Proposizione~\ref{prop:risparmio-proiezione} sulle sei matrici si ottiene
\[
s_{\text{strato}} \;=\; 1-\frac{3177}{16384} - 1{,}5\,\kappa - \frac{1{,}5}{d}
\;\approx\; 80{,}61\% - 1{,}5\,\kappa .
\]
Il coefficiente $1{,}5$ è la media, pesata per le dimensioni, del costo del basso rango nelle due famiglie di matrici: $2\kappa$ per ognuna delle quattro matrici $d\times d$ e $1{,}25\,\kappa$ per ognuna delle due matrici del feed-forward, cioè $(4\cdot 2\kappa\cdot d^2+2\cdot1{,}25\,\kappa\cdot4d^2)/12d^2=1{,}5\,\kappa$. Il termine $1{,}5/d$, dovuto alle unità aggiunte da scale e guadagni, vale lo $0{,}2\%$ per $d=768$ e tende a zero.

Ne seguono due comportamenti al crescere del modello:
\begin{itemize}
\item \textbf{con il rango proporzionale alla dimensione} ($\kappa$ fissato) il risparmio per strato è costante: con $\kappa=1/8$, come in STERA 30M ($96/768$), vale $61{,}9\%$ per ogni $d$;
\item \textbf{con il rango fissato} il termine del basso rango decresce come $1/d$ e il risparmio tende all'$80{,}61\%$ del formato.
\end{itemize}

\paragraph{Il modello intero.} Il risparmio garantito dai blocchi RowMask23 riguarda la parte strutturata delle matrici. Il risparmio di un modello STERA si ottiene invece contando tutte le componenti delle sue proiezioni e i parametri che non vengono compressi.

\begin{proposizione}[Risparmio di un modello STERA]
\label{prop:risparmio-modello}
Si consideri un Transformer decoder-only con $L$ strati, dimensione $d$, feed-forward $f$, vocabolario $V$ arrotondato a un multiplo di 128, tabella di $T$ posizioni ed embedding condiviso con la testa di uscita, in cui ogni proiezione è una proiezione STERA di rango $r$ (Definizione~\ref{def:proiezione-stera}). Siano
\[
\begin{aligned}
M &= V d + L\,(4d^2 + 2fd) &&\text{(valori delle proiezioni nella forma densa)},\\
R &= (V+d) + L\,\big(8d + 2(f+d)\big) &&\text{(somma di righe e colonne delle proiezioni)},\\
A &= T d + 2d\,(2L+1) + L\,(2d+f) &&\text{(posizioni, normalizzazioni e bias, non compressi)}.
\end{aligned}
\]
La forma densa equivalente conta $D = M + A$ valori, il modello STERA ne memorizza
\[
P \;=\; \underbrace{\frac{3177}{16384}\,M}_{\text{coordinate e scale}} \;+\; \underbrace{R}_{\text{guadagni}} \;+\; \underbrace{r\,R}_{\text{basso rango}} \;+\; \underbrace{A}_{\text{non compressi}},
\]
e il risparmio è
\[
S \;=\; 1 - \frac{P}{D} \;=\; \frac{M}{M+A}\left(1-\frac{3177}{16384}\right) \;-\; \frac{(1+r)\,R}{M+A}.
\]
\end{proposizione}

\begin{proof}
Si somma la Proposizione~\ref{prop:risparmio-proiezione} sulle proiezioni del modello: l'embedding ($V\times d$) e, in ogni strato, quattro matrici $d\times d$ e due matrici $f\times d$ e $d\times f$. La somma dei prodotti $m_i n_i$ è $M$ e la somma di $m_i+n_i$ è $R$. Ai valori delle proiezioni si aggiungono, in entrambe le forme, gli $A$ parametri non compressi.
\end{proof}

La seconda forma di $S$ rende esplicito il ruolo di ogni componente: il primo termine è l'$80{,}61\%$ garantito da RowMask23, ridotto della quota $A/(M+A)$ di parametri che non sono matrici; il secondo è il costo di guadagni e basso rango, l'unico che dipende dal rango $r$. Per STERA 30M ($L=6$, $d=768$, $f=3072$, $V=\num{50304}$, $T=\num{2048}$, $r=96$) si ha $M=\num{81100800}$, $R=\num{134016}$, $A=\num{1620480}$ e quindi
\[
P = \num{15726150} + \num{134016} + \num{12865536} + \num{1620480} = \num{30346182},
\qquad D = \num{82721280},
\qquad S = 63{,}32\%.
\]

\begin{table}[htbp]
\centering
\small
\begin{tabular}{lrr}
\toprule
Componente di STERA 30M & Valori & Quota \\
\midrule
Coordinate RowMask23 (\num{4950} blocchi) & \num{15721200} & $51{,}8\%$ \\
Scale dei blocchi & \num{4950} & $0{,}02\%$ \\
Guadagni per riga e colonna & \num{134016} & $0{,}4\%$ \\
Termine a basso rango ($r=96$) & \num{12865536} & $42{,}4\%$ \\
Posizioni, normalizzazioni, bias & \num{1620480} & $5{,}3\%$ \\
\midrule
Totale & \num{30346182} & \\
Forma densa equivalente & \num{82721280} & \\
Risparmio & $63{,}32\%$ & \\
\bottomrule
\end{tabular}
\caption{Valori memorizzati da STERA 30M per componente. Tutti i valori sono memorizzati e valutati in FP32.}
\label{tab:componenti-stera30m}
\end{table}

\begin{table}[htbp]
\centering
\small
\begin{tabular}{rrrrr}
\toprule
$d$ & senza basso rango & $r=96$ & $r=d/8$ & strato con $r=d/8$ \\
\midrule
512 & $78{,}5\%$ & $56{,}2\%$ & $63{,}6\%$ & $61{,}6\%$ \\
768 & $78{,}9\%$ & $63{,}3\%$ & $63{,}3\%$ & $61{,}7\%$ \\
1024 & $79{,}1\%$ & $67{,}1\%$ & $63{,}1\%$ & $61{,}7\%$ \\
2048 & $79{,}7\%$ & $73{,}3\%$ & $62{,}6\%$ & $61{,}8\%$ \\
4096 & $80{,}1\%$ & $76{,}7\%$ & $62{,}2\%$ & $61{,}8\%$ \\
\bottomrule
\end{tabular}
\caption{Risparmio rispetto alla forma densa di modelli con sei strati, vocabolario di \num{50304} token e \num{2048} posizioni, al variare della dimensione $d$ e del rango. Le prime tre colonne si riferiscono al modello intero, l'ultima al solo strato. Per $d=768$ e $r=96$ si ritrova STERA 30M.}
\label{tab:risparmio-crescita}
\end{table}

Con il rango proporzionale il risparmio del modello scende leggermente al crescere di $d$ perché il peso dell'embedding, che risparmia più degli strati, diminuisce, e tende al valore per strato del $61{,}9\%$. Con $r=96$ fissato il risparmio cresce con il modello e a $d=4096$ supera il $76\%$.
